\section{Method}
\label{sec:method}

\subsection{Memory entries and the store}
A memory entry $e$ is a self-contained question--answer pair with a
provenance set and a scalar \emph{energy}:
$e = (q_e, a_e, \mathcal{S}_e, E_e)$, where $\mathcal{S}_e$ is the set of
source documents the entry derives from. Following \citet{quek2026memo},
entries are produced by a reflection-style encoder that extracts explicit
and inferred facts, surfaces entities in both directions, and verifies
self-containment, so each $a_e$ is answerable without access to source
documents. A new entry spawns with energy $E_0 = 1.0$; energy is capped at
$\Emax = 5.0$; an entry is \emph{alive} iff $E_e > \varepsilon$ with
$\varepsilon = 10^{-9}$.

\paragraph{Retrieval is energy-blind.}
Given a query, the store ranks entries by a pure relevance score
$\rho(\text{query}, e)$ that \emph{never reads energy}; energy breaks ties
only. The invariant is load-bearing: if retrieval preferred high-energy
incumbents, selection pressure would come from the retriever rather than
from outcomes, and a wrong-but-entrenched entry could defend itself. The
default $\rho$ is a stdlib smoothed-IDF token overlap with a coverage
floor ($m_{\text{cov}} = 0.25$ of the query's IDF mass); any
$\text{text}\!\to\!\mathbb{R}^d$ embedder is a drop-in alternative under
the same invariant.

\subsection{Selection dynamics}
Let one \emph{cycle} present a batch of tasks drawn from a procedurally
regenerated environment, so no entry can overfit a fixed instance. For a
task, the query protocol returns an answer with a \emph{deciding} entry
$d$ (possibly $\varnothing$) and a set of \emph{supporting} entries
$\mathcal{P}$; the environment executes the implied action and returns a
reported outcome delta $\Delta \in \mathbb{R}$ (bytes freed, tests
gained). Three forces then act on energy.

\paragraph{(1) Credit along provenance.}
The measured delta is normalized and bounded before it becomes energy:
\begin{equation}
  c(\Delta) \;=\; \credit \cdot \tanh\!\left(\frac{\Delta}{\resscale}\right),
  \label{eq:credit}
\end{equation}
where $\credit = 0.6$ is the credit gain and $\resscale$ is the
environment's resource scale (e.g.\ $10^5$ bytes for the storage
environment). The $\tanh$ is the forgiveness mechanism: it keeps one
enormous outcome from making an entry immortal and one disaster from
instantly executing an entry that has been right many times. Credit is
distributed along provenance,
\begin{equation}
  \begin{aligned}
    \Delta E_d &= w_d\, c(\Delta), \quad
      w_d = \begin{cases} s & d \text{ juvenile}\\
                           1 & \text{otherwise}\end{cases}\\[2pt]
    \Delta E_j &= s\, c(\Delta)\ \ \forall j \in \mathcal{P},
  \end{aligned}
  \label{eq:provenance}
\end{equation}
with supporting share $s = 0.25$. When no single entry decides
(synthesized answers), credit spreads evenly, $\Delta E_j = c(\Delta)/|\mathcal{P}|$,
with probationary imports capped at the supporting share. Energy is then
updated with an upper clip,
\begin{equation}
  E_i \leftarrow \min\!\big(\Emax,\; E_i + \Delta E_i\big),
  \label{eq:apply}
\end{equation}
the lower side left unbounded so that death is decided at upkeep, not
here. The attribution rule associates a task outcome with consulted entries;
it is not a causal estimator. Test outcomes depend on the evaluation
and its coverage, and their reports remain vulnerable to corruption.

\paragraph{(2) Upkeep.}
Existence costs energy. Every cycle, each alive entry pays a fixed
upkeep, and any entry that thereby crosses the floor is buried:
\begin{equation}
  \begin{aligned}
    &E_i \leftarrow E_i - \upkeep \quad \forall i \ \text{alive};\\[2pt]
    &\text{bury } i \ \text{ iff } E_i \le \varepsilon
      \ \wedge\ i \notin (\text{protected} \cup \text{pinned}),
  \end{aligned}
  \label{eq:upkeep}
\end{equation}
with $\upkeep = 0.05$. Upkeep is what makes never-consulted entries
\emph{starve}: with no income they decay linearly to death. Pinned
entries pay upkeep but floor at zero instead of dying (rare-but-critical
knowledge whose payoff cadence is longer than the starvation horizon);
entries with an unsettled outcome are escrowed in \emph{protected} until
their verdict lands.

\paragraph{(3) Consolidation.}
Every $C = 5$ cycles, entries whose pairwise relevance similarity exceeds
$\tau = 0.55$ are merged. The merged entry unions provenance,
$\mathcal{S} = \bigcup_{i}\mathcal{S}_i$, and \emph{pools} energy under
the cap,
\begin{equation}
  E_{\text{merged}} = \min\!\Big(\Emax,\ \textstyle\sum_{i \in \text{cluster}} E_i\Big),
  \label{eq:merge}
\end{equation}
so capability per entry rises while the population shrinks. Consolidation
is hygiene, not safety: our ablations show outcomes are invariant to it.

\subsection{The cycle}
Algorithm~\ref{alg:cycle} states one cycle. The ordering matters: credit
is assigned per task as outcomes arrive, upkeep is charged once after all
tasks, and consolidation runs last and only periodically.

\begin{algorithm}[t]
\caption{One survival cycle}
\label{alg:cycle}
\begin{algorithmic}[1]
\Require store of entries; environment $\mathrm{Env}$; cycle index $t$
\For{each task $\tau$ in $\mathrm{Env}.\textsc{tasks}(t)$}
  \State $(\text{ans}, d, \mathcal{P}) \gets \textsc{Answer}(\text{store}, \tau)$
    \Comment{retrieval is energy-blind}
  \State $\Delta \gets \mathrm{Env}.\textsc{verify}(\tau, \text{ans})$
    \Comment{reported outcome delta}
  \State $c \gets \credit \cdot \tanh(\Delta / \resscale)$
    \Comment{Eq.~\ref{eq:credit}}
  \State apply $\Delta E_d, \Delta E_{j\in\mathcal{P}}$ via Eq.~\ref{eq:provenance}--\ref{eq:apply};
    advance trust lifecycle
\EndFor
\State \textbf{optionally} write the cycle's best positive trajectory as a
  new entry \Comment{must then survive on its own}
\State \textbf{for all} alive $i$: $E_i \gets E_i - \upkeep$;
  bury starved/executed entries \Comment{Eq.~\ref{eq:upkeep}}
\If{$t+1 \equiv 0 \pmod{C}$}
  \State merge clusters with similarity $\ge \tau$ \Comment{Eq.~\ref{eq:merge}}
\EndIf
\end{algorithmic}
\end{algorithm}

\subsection{Why poison dies without a label}
No rule in Algorithm~\ref{alg:cycle} refers to correctness, safety or
quality; the three death modes are emergent. \emph{Starvation}: an entry
never retrieved earns nothing and decays at rate $\upkeep$ per cycle.
\emph{Execution}: a poisoned entry that wins a query drives a harmful
action. In the filesystem-backed storage simulation, deleting a protected
file recreates that file and reports a modeled penalty of $3\times$ its
size; no restoration scratch space or physical I/O cost is measured.
That negative $c(\Delta)$
flows back along provenance until the entry crosses the floor.
\emph{Merging}: near-duplicate survivors pool into one. Bounded credit provides
a finite buffer against false reports (Eq.~\ref{eq:credit}). Ignoring upkeep,
an entry near $\Emax$ needs roughly $\Emax/\credit \approx 8$ maximal false
negatives to reach the floor; upkeep reduces that tolerance. Positive reported
outcomes can replenish the buffer. This is not a robustness guarantee: a
sufficiently large corruption budget defeats it, and a forgiveness counter can
also reset after positive outcomes. The experiments compare those boundaries.

\subsection{Trust lifecycle (optional)}
For adversarial imports, entries may enter on \emph{probation}
(they ride along and earn but may not decide until they have paid down a
number of net-positive settlements) or in a \emph{juvenile} admission
window (a locally minted entry earns at the capped supporting share, and
a single negative \emph{deciding} outcome denies admission outright). Both
default off and recover the base dynamics exactly; they are detailed in
the threat model and used only where stated.

\subsection{Who can corrupt outcome reports}
\label{sec:whocanlie}

We separate the ground-truth outcome from the report delivered to the
curator. The attacker changes only the delivered feedback within a stated
budget. The unchanged truth record scores useful-memory loss and harmful
decisions. A JUnit report, process exit code, or metrics export is software
output with its own access controls; calling it measured does not authenticate
it. Direct control of the store or evaluator is outside this constrained
feedback-corruption model and permits stronger attacks.

\paragraph{The historical CI lesson store.} The earlier integration for our target
(\S\ref{sec:regime}) settles a coding agent's lessons against test
outcomes: a merged pull request's pass-count delta pays the lessons that
advised it. Three parties write to that channel without any of them
controlling whether the tests truly pass---the pull request being
measured commits the store, the pull-request body names the tickets
settlement scrapes, and a contributor who adds a test adds a writer to
the channel. The middle one is the sharpest: open ticket ids are readable
in the committed store, so absent a check a contributor can paste
somebody else's in-flight ticket id and settle a stranger's decision at a
delta whose sign they choose. Our implementation refuses any id already
pending in the base commit for exactly this reason, and reports
\texttt{ticket\_provenance: unverified} when given no base to check
against---a defence that exists because the attack is reachable by
opening a pull request. The development GitHub/pytest profile instead binds a
ticket to repository, task, base commit, fixed evaluation, and run identity;
added tests earn no credit, and the recommended MCP configuration reserves
settlement for the host. These controls reduce accidental misattribution but
do not protect a store or evaluator that an attacker can rewrite directly.
\pointer{\S\ref{sec:cichannel} works all three through.}
\inplace{method-ci-channel}

\paragraph{Access assumptions.} The historical routes have different access
requirements: ticket substitution requires control of accepted pull-request
metadata, while evaluation changes require accepted code changes. Opening a
pull request alone does not guarantee either route reaches settlement.
Our controlled threat model grants alteration of delivered feedback under a
specified budget, without rewriting stored lessons. Unlike attacks that require
inserting lesson content (\S\ref{sec:related}), this attacker needs to
influence a number about work that genuinely happened. The adversary in
our experiments is weaker still---it observes the sign of the true delta
and nothing else. The same separation holds wherever settlement is
mediated: a storage agent reads a filesystem report rather than counting
sectors, a spend-cap agent reads a billing export that arrives a day
late. The limit is honest---where the settling process is the same
trusted component that performs the work and no third party can write to
the channel, this threat model does not apply and the attack budget is
zero, which is partly why we measure a curve from $b{=}0$ upward.

\subsection{Reporting rules}
\label{sec:reporting}

Every metric is computed from the resource movement that actually
occurred, never from what a lying measurement reported: scoring on the
reported number would credit an arm for being deceived in a convenient
direction. Two rules carry the paper's argument and are worth stating
here. \emph{Kill rate} is the fraction of seeds after which no surviving
entry from the poisoned source still advises an action---a claim about
advice that can be acted on, not about text remaining in the store.
\emph{Poison alive} counts surviving entries carrying the poisoned
provenance, inert ones included, and is reported beside the kill rate and
never instead of it, because a mechanism that removes nothing inert reads
$1.00$ and looks immune. Three readings in this paper were flattered
exactly that way before this count existed.
\pointer{\S\ref{sec:reporting-full} defines every reported quantity, and
a test pins that list against the code producing the tables.}
\inplace{method-reporting}
